% 第 11 题三视图：按固定源 PDF 矢量路径换算（1 pt → 0.03515 cm）。
% 原点取正视图左下角（源页 (524.5664, 375.2969) pt），y 向上。网格 28.1953 pt = 2 cm。
% 正视图：2×2 格、左上格缺（中竖贯通、中横仅左半）；侧视图：2×2 格、右下格缺（中横贯通、
% 中竖下半为虚线——源内线型如此，按源照录）；俯视图：左列满深、右列只余上半。
% 尺寸标注：横向 4 条 + 竖向 2 条尺寸线（均两段让位给数字），6 条竖直刻度线、3 条水平刻度线，
% 12 个实心箭头（横向一律朝外）；「2」「正视图」「侧视图」「俯视图」共 9 个标签。
\begin{tikzpicture}[line cap=round,line join=round,font=\small,x=0.03515cm,y=0.03515cm]
% ===== 正视图（左格）=====
\draw (0,0)--(56.39,0)--(56.39,56.39)--(28.20,56.39)--(28.20,28.20)--(0,28.20)--cycle;
\draw (28.20,0)--(28.20,28.20);                 % 中竖下半（实线）
% ===== 侧视图（右格，右移 84.58）=====
\draw (84.58,0)--(140.97,0)--(140.97,28.20)--(112.78,28.20)--(112.78,56.39)--(84.58,56.39)--cycle;
\draw (84.58,28.20)--(112.78,28.20);            % 中横左半（实线）
\draw[dash pattern=on 2.99pt off 2.99pt] (112.78,0)--(112.78,28.20);  % 中竖下半（虚线，照源）
% ===== 俯视图 =====
\draw (0,-28.20)--(28.20,-28.20)--(28.20,-84.58)--(0,-84.58)--cycle;   % 左列满深
\draw (28.20,-28.20)--(56.39,-28.20)--(56.39,-56.39)--(28.20,-56.39);  % 右列只余上半（第四边与左列右边重合）
% ===== 尺寸线（横向 4 条 @y=62.03，两段让位）=====
\draw (3.59,62.03)--(24.61,62.03);   \draw (31.78,62.03)--(52.80,62.03);     % 正视图两格
\draw (88.17,62.03)--(109.19,62.03); \draw (116.36,62.03)--(137.39,62.03);  % 侧视图两格
% ===== 尺寸线（竖向 2 条 @x=62.03，两段让位）=====
\draw (62.03,31.78)--(62.03,52.81);  % 上段（高 2）
\draw (62.03,3.59)--(62.03,24.61);   % 下段（高 2）
% ===== 竖直刻度线（6 条，@y 59.21–67.67）=====
\draw (0,59.21)--(0,67.67);        \draw (28.20,59.21)--(28.20,67.67);
\draw (56.39,59.21)--(56.39,67.67); \draw (84.58,59.21)--(84.58,67.67);
\draw (112.78,59.21)--(112.78,67.67); \draw (140.97,59.21)--(140.97,67.67);
% ===== 水平刻度线（3 条，@x 59.21–67.67）=====
\draw (59.21,0)--(67.67,0);
\draw (59.21,28.20)--(67.67,28.20);
\draw (59.21,56.39)--(67.67,56.39);
% ===== 实心箭头（12 个，底 2.98 长 3.98，横向一律朝外）=====
\fill (0,62.03)--(3.98,60.54)--(3.98,63.52)--cycle;         % 正左格：左端←
\fill (28.20,62.03)--(24.22,60.54)--(24.22,63.52)--cycle;   % 正左格：右端→
\fill (28.20,62.03)--(32.18,60.54)--(32.18,63.52)--cycle;   % 正右格：左端←
\fill (56.39,62.03)--(52.41,60.54)--(52.41,63.52)--cycle;   % 正右格：右端→
\fill (84.58,62.03)--(88.56,60.54)--(88.56,63.52)--cycle;   % 侧左格：左端←
\fill (112.78,62.03)--(108.80,60.54)--(108.80,63.52)--cycle;% 侧左格：右端→
\fill (112.78,62.03)--(116.76,60.54)--(116.76,63.52)--cycle;% 侧右格：左端←
\fill (140.97,62.03)--(136.99,60.54)--(136.99,63.52)--cycle;% 侧右格：右端→
\fill (62.03,56.39)--(60.54,52.41)--(63.52,52.41)--cycle;   % 上段：上端↑
\fill (62.03,28.20)--(60.54,32.18)--(63.52,32.18)--cycle;   % 上段：下端↓
\fill (62.03,28.20)--(60.54,24.22)--(63.52,24.22)--cycle;   % 下段：上端↑
\fill (62.03,0)--(60.54,3.99)--(63.52,3.99)--cycle;         % 下段：下端↓
% ===== 标签（照源字形框中心，显式坐标）=====
\node[anchor=center,inner sep=0] at (14.30,66.90) {\(2\)};
\node[anchor=center,inner sep=0] at (42.51,66.90) {\(2\)};
\node[anchor=center,inner sep=0] at (98.92,66.90) {\(2\)};
\node[anchor=center,inner sep=0] at (127.12,66.90) {\(2\)};
\node[anchor=center,inner sep=0] at (67.89,41.51) {\(2\)};
\node[anchor=center,inner sep=0] at (67.89,13.31) {\(2\)};
\node[anchor=center,inner sep=0] at (28.40,-13.54) {正视图};
\node[anchor=center,inner sep=0] at (113.01,-13.54) {侧视图};
\node[anchor=center,inner sep=0] at (28.40,-98.09) {俯视图};
\end{tikzpicture}
